% Part: first-order-logic
% Chapter: model-theory
% Section: partial-iso

\documentclass[../../../include/open-logic-section]{subfiles}

\begin{document}

\olfileid{mod}{bas}{pis}
\section{جزوی ہم شکلیاں}

\begin{defn}
  دو !!{structure}s $\Struct{M}$ تے $\Struct{N}$ دتیاں ہون۔
  $\Struct{M}$ توں $\Struct{N}$ ول اک \emph{جزوی ہم شکلی}
  (partial isomorphism) اک محدود جزوی فنکشن $p$ اے، جیہڑا
  $\Domain M$ دے عناصر نوں دلیل بنا کے $\Domain N$ وچ قدراں دیندا
  اے، تے اپنے دائرے اُتے \olref[iso]{defn:isomorphism} دیاں
  ہم شکلی والیاں شرطاں پوری کردا اے:
  \begin{enumerate}
  \item $p$ !!{injective} اے؛
  \item ہر !!{constant}~$c$ لئی: جے $p(\Assign{c}{M})$ دی قدر
    متعین ہووے، تاں $p(\Assign{c}{M}) = \Assign{c}{N}$؛
  \item ہر $n$-تھانواں والے !!{predicate} $P$ لئی: جے $a_1$، \dots، $a_n$
    سارے $p$ دے دائرے وچ ہوون، تاں $\langle a_1, \dots, a_n\rangle \in
    \Assign P M$ اُدوں تے صرف اُدوں جدوں $\langle p(a_1), \dots, p(a_n) \rangle
    \in \Assign P N$؛
  \item ہر $n$-تھانواں والے !!{function} $f$ لئی: جے $a_1$، \dots، $a_n$
    سارے $p$ دے دائرے وچ ہوون، تاں $p(\Assign f M (a_1, \dots,a_n))
    = \Assign f N (p(a_1), \dots, p(a_n))$۔
  \end{enumerate}
  $p$ دے محدود ہون دا مطلب ایہ اے کہ $\dom{p}$ محدود اے۔
  دوہاں ساختاں دی زبان اکّو اے۔ فنکشن والی شرط وچ، دائرے دے عناصر
  اُتے فنکشن دی قدر وی جزوی نقشے دے دائرے وچ ہونی چاہیدی اے، نہیں تاں
  کھبے پاسے دی قدر متعین نہیں ہوئے گی۔ پورے ہدف دائرے اُتے ہون دی شرط نہیں۔
\end{defn}

دھیان دیو کہ خالی فنکشن~$\emptyset$ ہر دو !!{structure}s وچکار
ہمیشہ اک جزوی ہم شکلی اے۔
ایتھے کتاب دی زبان مطابق ہر محمول تے فنکشن دے تھانواں دی تعداد مثبت اے؛
مستقل دی شرط وی صرف اوہدوں لگدی اے جدوں اوہدی قدر نقشے دے دائرے وچ ہووے۔

\begin{defn}\ollabel{defn:partialisom}
  دو !!{structure}s $\Struct{M}$ تے $\Struct{N}$
  \emph{جزوی طور اُتے ہم شکل} نیں، جس نوں $\Struct{M} \iso[p]
  \Struct{N}$ لکھدے آں، اُدوں تے صرف اُدوں جدوں جزوی ہم شکلیاں
  دا اک غیر خالی سیٹ $I$ ہووے، جیہڑیاں $\Struct{M}$ تے $\Struct{N}$
  وچکار ہون، تے ایہ سیٹ \emph{اگے-پیچھے} (back-and-forth) خاصیت پوری کرے:
  \begin{enumerate}
  \item (\emph{اگے}) ہر $p \in I$ تے $a \in \Domain M$ لئی
    کوئی $q \in I$ اے جیہدے لئی $p \subseteq q$، تے $a$،
    $q$ دے دائرے وچ اے؛
  \item (\emph{پیچھے}) ہر $p \in I$ تے $b \in \Domain N$ لئی
    کوئی $q \in I$ اے جیہدے لئی $p \subseteq q$، تے $b$،
    $q$ دیاں حاصل قدراں وچ اے۔
  \end{enumerate}
\end{defn}

\begin{thm}\ollabel{thm:p-isom1}
  جے $\Struct{M} \iso[p] \Struct{N}$ تے $\Struct{M}$ تے
  $\Struct{N}$ !!{enumerable} ہون، تاں $\Struct{M} \iso
  \Struct{N}$۔
\end{thm}

\begin{proof}
  کیوں جے $\Struct{M}$ تے $\Struct{N}$ !!{enumerable} نیں، فرض کرو
  $\Domain{M} =
  \{a_0, a_1, \ldots \}$ تے $\Domain{N} = \{b_0, b_1, \ldots \}$۔
  جے کوئی دائرہ محدود ہووے، تاں سارے عنصر لسٹ کرکے آخری عنصر نوں دہراؤ؛
  کتاب مطابق دائرہ غیر خالی اے۔ لسٹ دے ہر اصل عنصر نوں ہن وی اک مرحلہ ملدا اے۔
  کسے وی $p_0 \in I$ توں شروع کرکے اسیں جزوی ہم شکلیاں دا
  اک ودھدا سلسلہ $p_0 \subseteq p_1 \subseteq p_2
  \subseteq \cdots$ ایویں مقرر کردے آں:
  \begin{enumerate}
  \item جے $n+1$ طاق ہووے، یعنی $n = 2r$، تاں اگے والی خاصیت
    نال کوئی $p_{n+1} \in I$ لبھو جیہدے لئی $p_n \subseteq p_{n+1}$
    تے $a_r$، $p_{n+1}$ دے دائرے وچ ہووے؛
  \item جے $n+1$ جفت ہووے، یعنی $n =2r+1$، تاں پیچھے والی خاصیت
    نال کوئی $p_{n+1} \in I$ لبھو جیہدے لئی $p_n \subseteq p_{n+1}$
    تے $b_r$، $p_{n+1}$ دیاں حاصل قدراں وچ ہووے۔
  \end{enumerate}
ایس طرح پہلی اگلی چال پہلی لسٹ دے شروع والے عنصر نوں، تے پہلی پچھلی چال
دوجی لسٹ دے شروع والے عنصر نوں شامل کردی اے؛ کوئی شروع والا عنصر چھٹدا نہیں۔
ہن جے اسیں ایہ رکھئیے:
\[
p = \bigcup_{n\ge 0} p_n,
\]
تاں $p$، $\Struct{M}$ تے $\Struct{N}$ وچکار اک ہم شکلی اے۔
\end{proof}

\begin{prob}
  پوری تفصیل نال دکھاؤ کہ
  \olref[mod][bas][pis]{thm:p-isom1} وچ مقرر کیتا $p$ واقعی اک
  ہم شکلی اے۔
\end{prob}

\begin{thm}\ollabel{thm:p-isom2}
  فرض کرو $\Struct{M}$ تے $\Struct{N}$ اک خالص تعلقی
  !!{language} دیاں !!{structure}s نیں (یعنی !!a{language} جیہدے
  وچ صرف !!{predicate}s ہون، کوئی !!{function}s یا مستقل نہ ہون)۔
  تاں جے $\Struct{M} \iso[p] \Struct{N}$، ایہ وی ہووے گا کہ
  $\Struct{M} \elemequiv
  \Struct{N}$۔
\end{thm}

\begin{proof}
  ایس دعوے وچ جزوی ہم شکلی اوپر دتے اگے-پیچھے والے غیر خالی سیٹ وچوں
  ہونی چاہیدی اے۔ فارمولے دے سارے آزاد متغیر تھلے دتی متغیر لسٹ وچ ہون۔
  !!{formula}s اُتے استقرا کرکے دکھایا جاندا اے کہ جے $a_1$، \dots، $a_n$ تے
  $b_1$، \dots، $b_n$ ایسے ہون کہ اگے-پیچھے والے سیٹ دی اک جزوی ہم شکلی $p$ ہر
  $a_i$ نوں $b_i$ ول بھیجے تے $s_1(x_i) =a_i$ تے $s_2(x_i) =b_i$
  ($i =1$، \dots،~$n$ لئی)، تاں $\Sat{M}{!A}[s_1]$ اُدوں تے
  صرف اُدوں جدوں $\Sat{N}{!A}[s_2]$۔ $n=0$ دا حال
  $\Struct{M} \elemequiv \Struct{N}$ دیندا اے۔
  اکیلی جزوی ہم شکلی، جیہڑی ایس توسیع والے سیٹ دا حصہ نہ ہووے، کافی نہیں۔
\end{proof}

\begin{rem}
جے !!{function}s وی ہون، تاں پچھلا نتیجہ فیر وی درست اے، پر
$p$ توں پیدا ہون والی اوہ ہم شکلی ویکھنی پیندی اے جیہڑی
$\Struct{M}$ دی، $a_1$، \dots، $a_n$ توں پیدا ہوئی
ذیلی!!{structure} تے $\Struct{N}$ دی، $b_1$، \dots، $b_n$
توں پیدا ہوئی ذیلی!!{structure} وچکار ہووے۔
ایتھے وی اگے-پیچھے والا پورا مفروضہ قائم اے، صرف اکیلا جزوی نقشہ کافی نہیں۔
جملیاں دے خالی متغیر لسٹ والے حال وچ غیر خالی نقشے تک توسیع لے کے کم کیتا
جا سکدا اے؛ فارمولے وچ آؤن والیاں مستقل قدراں وی متعلقہ توسیع وچ شامل کرو۔
\end{rem}

پچھلے نتیجے نوں مرحلیاں وچ وی وکھایا جا سکدا اے: !!a{formula}
وچ اک دوجے دے اندر آئے مقدار نماواں دی تعداد تے متعلقہ جزوی
ہم شکلیاں نوں کنی واری ودھایا جا سکدا اے، انہاں وچ ربط قائم کرکے۔

\begin{defn}
  ہر !!{formula}~$!A$ لئی $!A$ دا \emph{مقدار نمائی درجہ}
  (quantifier rank)، جیہنوں $\QuantRank{!A} \in \Nat$ لکھدے آں،
  استقرائی طور اُتے $!A$ وچ اک دوجے دے اندر آئے مقدار نماواں دی
  سبھ توں وڈی تعداد دے طور اُتے مقرر ہوندا اے۔ دو
  !!{structure}s $\Struct{M}$ تے $\Struct{N}$ \emph{$n$-ہم ارز}
  نیں، جیہنوں $\Struct{M} \elemequiv[n] \Struct{N}$ لکھدے آں،
  جے اوہ مقدار نمائی درجہ~$n$ توں گھٹ یا اوہدے برابر رکھن والے
  سارے جملیاں بارے متفق ہون۔
\end{defn}

\begin{prop}\ollabel{prop:qr-finite}
  $\Lang{L}$ اک محدود خالص تعلقی !!{language} ہووے، یعنی ایسی
  !!{language} جیہدے وچ محدود گنتی دے !!{predicate}s تے
  !!{constant}s ہون، تے کوئی !!{function}s نہ ہوون۔ تاں ہر
  $n \in \Nat$ لئی !!{language} $\Lang{L}$ وچ پہلے درجے دے
  !!{sentence}s، جنہاں دا مقدار نمائی درجہ $n$ توں ودھ نہ ہووے،
  منطقی ہم ارزی تک صرف محدود گنتی وچ نیں۔
  ایس قضیے وچ مستقل نشانیاں دی محدود تعداد دی اجازت اے؛ پچھلے قضیے والی
  بنا مستقل والی قید ایتھے نہیں لائی گئی۔ جملیاں دے نال ایہ نتیجہ اوہناں
  فارمولیاں لئی وی درست اے جیہڑے اک مقرر محدود لسٹ دے متغیراں توں باہر
  کوئی آزاد متغیر نہیں رکھدے۔ محدود گنتی ہم ارزی جماعتاں دی اے، لفظی فارمولیاں دی نہیں۔
\end{prop}

\begin{proof}
  $n$ اُتے استقرا نال؛ فارمولیاں والے روپ وچ آزاد متغیراں دی محدود لسٹ مقرر رکھو۔
\end{proof}

\begin{defn}
  !!a{structure}~$\Struct{M}$ دتی ہووے؛ $\Domain M^{<\omega}$،
  $\Domain{M}$ دے عناصر اُتے بنے سارے محدود سلسلیاں دا سیٹ ہووے۔
  اسیں $\mathbf{a},
  \mathbf{b}, \mathbf{c}, \ldots$ نوں عناصر دے محدود سلسلیاں
  لئی ورتدے آں۔ جے $\mathbf{a} \in \Domain{M}^{<\omega}$ تے
  $a \in \Domain{M}$، تاں $\mathbf{a}a$، $\mathbf{a}$ نال
  $a$ دا \emph{جوڑ} (concatenation) ظاہر کردا اے۔
\end{defn}

\begin{defn}
  !!{structure}s $\Struct{M}$ تے $\Struct{N}$ دتیاں ہون۔
  اسیں برابر لمبائی والے سلسلیاں وچکار تعلقات
  $I_n \subseteq \Domain M^{<\omega} \times \Domain N^{<\omega}$
  نوں $n$ اُتے استقرا نال ایویں مقرر کردے آں:
   \begin{enumerate}
   \item $I_0(\mathbf{a},\mathbf{b})$ اُدوں تے صرف اُدوں جدوں
     $\mathbf{a}$ تے $\mathbf{b}$، $\Struct{M}$ تے
     $\Struct{N}$ وچ اکّو جوہری !!{formula}s پورے کردے ہون؛ یعنی
     جے $s_1(x_i) = a_i$ تے $s_2(x_i) =
     b_i$، تے $!A$ جوہری ہووے تے اوہدے سارے !!{variable}s،
     $x_1$، \dots،~$x_n$ وچوں ہون، تاں $\Sat{M}{!A}[s_1]$ اُدوں تے
     صرف اُدوں جدوں~$\Sat{N}{!A}[s_2]$۔
   \item $I_{n+1} (\mathbf{a},\mathbf{b})$ اُدوں تے صرف اُدوں جدوں ہر
     $a\in \Domain M$ لئی کوئی $b\in \Domain N$ ہووے جیہدے لئی $I_n
     (\mathbf{a}a,\mathbf{b}b)$، تے الٹی سمت وچ وی ایویں ہووے۔
   \end{enumerate}
  جوہری شرط وچ متغیر لسٹ دی آخری گنتی سلسلے دی لمبائی اے؛ اوہنوں
  تعلق دے استقرائی درجے دی گنتی نال نہ ملاؤ۔ دوہاں سلسلیاں دی لمبائی برابر اے۔
\end{defn}


\begin{defn}
  $\Struct{M} \approx_n \Struct{N}$ لکھو جے
  $I_n(\emptyseq,\emptyseq)$، $\Struct{M}$ تے
  $\Struct{N}$ لئی درست ہووے (جتھے $\emptyseq$ خالی سلسلہ اے)۔
\end{defn}

\begin{thm}\ollabel{thm:b-n-f}
  $\Lang{L}$ اک خالص تعلقی !!{language} ہووے۔ تاں $I_n
  (\mathbf{a},\mathbf{b})$ توں نکلدا اے کہ ہر $!A$، جیہدے لئی
  $\QuantRank{!A} \le n$، بارے $\Sat{M}{!A}[\mathbf{a}]$ اُدوں تے
  صرف اُدوں جدوں $\Sat{N}{!A}[\mathbf{b}]$ (ایتھے وی $\mathbf{a}$،
  $!A$ نوں پورا کردا اے جے کوئی وی $s$، جیہدے لئی
  $s(x_i) = a_i$، $!A$ نوں پورا کرے)۔ ہور ایہ کہ جے $\Lang{L}$
  محدود ہووے، تاں الٹا دعویٰ وی درست اے۔
  ایتھے ہر فارمولے دے سارے آزاد متغیر دتی متغیر لسٹ وچ ہون؛ ایس لئی
  لسٹ نال ملدیاں قدر بندیاں باقی متغیراں اُتے کیہ قدراں دیندیاں نیں، فرق نہیں پیندا۔
\end{thm}

\begin{proof}
  $I_n(\mathbf{a},\mathbf{b})$ توں ایہ نکلن دا ثبوت کہ
  $\mathbf{a}$ تے $\mathbf{b}$، مقدار نمائی درجہ $n$ توں
  ودھ نہ رکھن والے اکّو !!{formula}s پورے کردے نیں،
  $!A$ اُتے سادے استقرا نال ملدا اے۔ الٹے دعوے لئی اسیں
  $n$ اُتے استقرا کردے آں، تے \olref{prop:qr-finite} ورتدے آں؛
  ایہ یقینی بناؤندا اے کہ ہر $n$ لئی ایس مقدار نمائی درجے دے
  غیر ہم ارز !!{formula}s دی زیادہ توں زیادہ محدود گنتی اے۔
  ایتھے اوپر ثابت کیتا مقرر آزاد متغیر لسٹ والا روپ ورتیا گیا اے۔

  $n=0$ لئی ایہ مفروضہ کہ $\mathbf{a}$ تے $\mathbf{b}$ اکّو
  بے مقدار نما !!{formula}s پورے کردے نیں، ایہ وی دیندا اے کہ اوہ
  اکّو جوہری فارمولے پورے کردے نیں؛ ایس لئی
  $I_0(\mathbf{a},\mathbf{b})$۔

  $n+1$ دے حال وچ فرض کرو $\mathbf{a}$ تے $\mathbf{b}$،
  مقدار نمائی درجہ $n+1$ توں ودھ نہ رکھن والے اکّو
  !!{formula}s پورے کردے نیں۔ $I_{n+1}(\mathbf{a},\mathbf{b})$
  دکھاؤن لئی دونوں سمتاں دکھاؤنیاں نیں۔ اگلی سمت لئی ہر $a \in \Domain M$ لئی
  کوئی $b \in
  \Domain N$ ہووے جیہدے لئی $I_n(\mathbf{a}a,\mathbf{b}b)$؛
  تے استقرائی مفروضے نال ایہ دکھانا کافی اے کہ ہر $a \in
  \Domain M$ لئی کوئی $b \in \Domain N$ ہووے جیہدے لئی
  $\mathbf{a}a$ تے $\mathbf{b}b$ مقدار نمائی درجہ
  $n$ توں ودھ نہ رکھن والے اکّو !!{formula}s پورے کردے ہون۔

  کوئی $a \in \Domain M$ لو؛ $!T^a_n$ لئی اوہناں !!{formula}s
  $!B(x,\mathbf{y})$ دیاں منطقی ہم ارزی جماعتاں وچوں اک اک نمائندہ چنو
  جنہاں دا درجہ $n$ توں ودھ نہیں تے جنہاں نوں $\mathbf{a}a$، $\Struct{M}$ وچ پورا کردا اے۔
  $!T^a_n$ ہن چنے ہوئے سچ نمائندیاں دا محدود سیٹ اے، سارے لفظی فارمولیاں دا نہیں۔
  ایس لئی اوہناں ساریاں دی عطف نال پہلے درجے دا اک
  !!{formula} بنا لو؛ عطف مقدار نمائی درجہ نہیں ودھاؤندی۔ ایس توں نکلدا اے کہ
  $\mathbf{a}$، $\lexists[x][!T^a_n(x,\mathbf{y})]$ نوں پورا
  کردا اے، جیہدا مقدار نمائی درجہ $n+1$ توں ودھ نہیں۔
  مفروضے مطابق $\mathbf{b}$، $\Struct{N}$ وچ ایہو
  !!{formula} پورا کردا اے، ایس لئی کوئی $b \in \Domain N$ اے
  جیہدے نال $\mathbf{b}b$، $!T^a_n$ نوں پورا کردا اے؛ خاص طور
  اُتے $\mathbf{b}b$، مقدار نمائی درجہ $n$ توں ودھ نہ رکھن
  والے اوہو !!{formula}s پورے کردا اے جیہڑے $\mathbf{a}a$ پورے
  کردا اے۔ دونوں رخ ایس لئی ملدے نیں کہ اصل لسٹ لئی جھوٹے فارمولے
  دی نفی وی ایس درجے وچ سچی اے؛ اوہدی ہم ارزی جماعت دا نمائندہ وی
  چنے سیٹ وچ اے، اس لئی نئی لسٹ اوہ جھوٹا فارمولا سچ نہیں کر سکدی۔
  ایس طرح ہر $b \in \Domain N$ لئی کوئی
  $a\in \Domain M$ لبھدا اے جیہدے نال $\mathbf{a}a$ تے
  $\mathbf{b}b$ مقدار نمائی درجہ $n$ توں ودھ نہ رکھن والے
  اکّو !!{formula}s پورے کردے نیں؛ ایس نال ثبوت پورا ہوندا اے۔
\end{proof}

\begin{cor}\ollabel{cor:b-n-f}
  جے $\Struct{M}$ تے $\Struct{N}$ اک محدود !!{language} دیاں
  خالص تعلقی !!{structure}s نیں، تاں $\Struct{M} \approx_n\Struct{N}$
  اُدوں تے صرف اُدوں جدوں $\Struct{M} \elemequiv[n] \Struct{N}$۔
  خاص طور اُتے $\Struct{M} \elemequiv \Struct{N}$ اُدوں تے
  صرف اُدوں جدوں ہر $n$ لئی $\Struct{M} \approx_n \Struct{N}$۔
\end{cor}

\end{document}
